% Lecture visuelle directe avec la vision du modèle ; aucun OCR.
% Source : PDF original, vue 011, côté droit, page manuscrite 21, à partir du titre du chapitre second.
% Statut : transcrit ; relecture indépendante à effectuer.
\chapter{Gouvernement du Directoire exécutif. Élection du pape Pie VII}

Notre malheureux pays étoit bien loin d'être parvenu à un état parfait de bonheur ; toutefois, en comparant sa situation actuelle avec celle à laquelle il venoit d'échapper, on eût été presque tenté de la croire heureuse. La Convention nationale, cette assemblée de dégoûtante et sanglante mémoire, après avoir fait mourir son roi et avoir ensuite courbé sa tête coupable sous le joug hideux de Robespierre, termina enfin, quelque
% Source : PDF original, vue 012, côté gauche, page manuscrite 22.
temps après la chute de ce scélérat, sa trop longue carrière. Avant de se séparer elle nous dota d'une nouvelle constitution. Une commission de cinq membres, qu'on appeloit directeurs, formoit sous le nom de Directoire le pouvoir exécutif ; deux assemblées délibérantes, l'une nommée Conseil des Anciens parce qu'il falloit avoir au moins 40 ans pour en faire partie, et l'autre Conseil des Cinq-Cents à cause du nombre de ses membres, exerçoient le pouvoir législatif. Ces conseils nommoient les directeurs et devoient eux-mêmes être élus par les assemblées électorales des départements. Mais pour le moment, les conventionnels, pour conserver le plus longtemps possible un pouvoir dont ils eussent désiré ne jamais se départir, s'étoient arrogé le droit de les composer en grande partie des membres mêmes de la Convention, sauf à les renouveller ensuite d'année en année, par tiers, ainsi que le prescrivoit la Constitution.

Bientôt arriva le moment où le premier tiers devoit être remplacé. Les excès de la Convention avoient amené dans l'opinion publique une réaction bien prononcée ; elle se manifesta dans le choix de ce premier tiers. Son arrivée opéra un changement notable dans l'esprit des deux conseils, mais surtout dans celui du Conseil des Cinq-Cents. On osa plaider à la tribune, quoique avec un peu de mollesse, la cause de la religion. La société [lecture incertaine : sembloit alors revenir] à des idées d'ordre\footnote{La lecture de cette proposition reste incertaine ; voir le registre des questions, Q009. (Note du transcripteur).}, le mouvement qui y revenoit se faisoit sentir partout. À la vérité, les lois contre les catholiques subsistoient
% Source : PDF original, vue 012, côté droit, page manuscrite 23.
toujours ; mais dans la pratique on ne les exécutoit plus. Le gouvernement n'en protégeoit pas moins de tout son pouvoir l'Église constitutionnelle. La Terreur avoit dispersé et persécuté le clergé de cette Église schismatique tout comme le clergé catholique. Quelques-uns de ses membres durent à cette persécution commune le bonheur de reconnoître leurs fautes, et plusieurs obtinrent ainsi un glorieux martyre qui leur mérita leur expiation. Mais la plupart au contraire s'endurcirent davantage. Un abîme appelle un autre abîme, \textit{Abyssus abyssum invocat.}\footnote{Psaume 41 (42), 8 : « Un abîme appelle un autre abîme. » (Note du transcripteur).} C'étoit pour eux en grand nombre la crainte de se voir enlever tout moyen d'existence qui les avoit rendus schismatiques. La crainte de la mort en fit des apostats ; ils livrèrent lâchement leurs lettres de prêtrise ; il s'en trouva parmi eux qui furent assez déhontés pour se soumettre à annoncer publiquement au peuple que jusqu'ici leur bouche n'avoit proféré que des mensonges.

Aussitôt que nous fûmes sortis des jours de la Terreur, les débris de cette Église séparée cherchèrent à se réunir. Les constitutionnels parvinrent à se faire remettre en jouissance d'une quantité assez considérable d'églises paroissiales dans le diocèse de Lyon ; dans notre ville même ils en obtinrent quatre. Leur évêque intrus, M. Lamourette, avoit témoigné de repentir de ses fautes, et étoit mort victime de la Terreur ; ils le remplacèrent par un M. Primat qui établit son siège épiscopal dans l'église de St Nizier, car on n'avoit pas voulu leur donner celle de St Jean. L'Église constitutionnelle étoit donc réorganisée. Néanmoins, malgré ces avantages
% Source : PDF original, vue 013, côté gauche, page manuscrite 24.
extérieurs dont elle étoit en possession, le schisme comptoit à Lyon bien peu de partisans. Les habitants de cette ville étoient tout comme aujourd'hui partagés en deux classes : les vrais catholiques et les hommes sans religion. Ceux-là ne vouloient pas plus de l'Église séparée que de la véritable Église.

\section*{État du diocèse de Lyon}
% Mention marginale de la page manuscrite 24.
L'administration du diocèse ne laissoit rien à désirer. Je dois d'autant mieux vous en donner ici quelque idée que les moyens que nos supérieurs ecclésiastiques ont su employer avec tant de sagesse pour nous conserver la foi peuvent bien être considérés comme une des causes éloignées qui ont amené la formation de la Congrégation. Mgr de [nom incertain : Marbœuf] avoit été nommé archevêque de Lyon au commencement de la Révolution. Il ne lui a jamais été possible de prendre en personne possession de son siège ; mais il n'a jamais cessé jusqu'à sa mort d'administrer le diocèse du fond de sa retraite avec autant de zèle que de prudence. Parmi les prêtres qui, demeurés à Lyon, s'étoient dévoués à affronter toutes les rigueurs de la persécution pour ne pas laisser sans secours les âmes fidèles, il en avoit choisi douze qu'il avoit nommés vicaires généraux et dont il avoit formé un conseil, se réservant de diriger leur action par correspondance. À la tête de ce conseil se distinguoit par ses talents administratifs et sa courageuse fermeté M. Linsolas, ancien prêtre habitué de l'église de St Nizier, mort en 1828 chanoine de St Jean et vicaire général d'honneur. La dette que le diocèse de Lyon a contractée envers ce vénérable ecclésiastique est immense ; heureusement Dieu est assez riche pour l'acquitter.

% Source : PDF original, vue 013, côté droit, page manuscrite 25.
M. Linsolas présenta au conseil un plan d'administration qui fut adopté dans toutes ses parties. Le diocèse fut divisé en cercles de mission formés par la réunion de plusieurs paroisses ; on attacha à chacun de ces cercles un chef de mission et plusieurs missionnaires ; de cette manière ces prêtres, n'ayant point de résidence fixe, mais transportant successivement leur domicile suivant que l'exigeoient les circonstances, dans l'une ou l'autre des habitations des paroisses qui composoient la mission qui leur étoit confiée, étoient beaucoup moins exposés à tomber entre les mains des persécuteurs, et si on en excepte toutefois les jours les plus affreux de la Terreur, les fidèles purent toujours se procurer avec facilité les secours de la religion. Seulement lorsqu'il y avoit recrudescence de persécution, on redoubloit de prudence. On avoit choisi dans chaque village des laïcs recommandables pour exercer comme dans les missions des sauvages les fonctions de catéchistes et présider à l'assemblée des fidèles en l'absence du prêtre. Le nombre des maisons dans lesquelles se célébroient à Lyon les saints mystères a toujours été considérable ; il s'est élevé progressivement jusqu'à celui de 400. À l'époque dont nous nous occupons on jouissoit d'une espèce de liberté ; ainsi le service divin s'exerçoit à peu près publiquement dans un grand nombre de ces oratoires. Dans les campagnes, les catholiques se réunissoient presque partout dans de vastes granges qu'on transformoit en oratoires. Plusieurs églises qui n'étoient pas occupées par les constitutionnels furent même réparées par les catholiques dans les paroisses
% Source : PDF original, vue 014, côté gauche, page manuscrite 26.
où ils se trouvoient en très grande majorité. Les vicaires généraux essayèrent de profiter de ce moment de calme pour assurer la perpétuité du sacerdoce en ouvrant un grand séminaire. Ce fut vers le commencement de 1797 qu'ils y réunirent les jeunes gens qui se destinoient à l'état ecclésiastique. Le séminaire étoit placé à la montée du Garillan, dans une maison occupée autrefois par une Providence et où sont maintenant les dames Carmélites.

La tendance générale de l'opinion paroissoit alors, ainsi que nous l'avons déjà remarqué, de plus en plus favorable à une heureuse solution des affaires de la religion et de celles de la patrie. Suivant toutes les apparences, l'arrivée d'un second tiers dans les conseils devoit donner aux ennemis du désordre et des révolutions une majorité prononcée. Dans cette confiance, non seulement les réfugiés qu'une loi avoit rappelés en France s'étoient empressés d'y rentrer, mais les émigrés eux-mêmes, quoique toujours proscrits, commençoient à reparoître, et les autorités de la République, dans l'attente des événements, fermoient les yeux. Hélas ! il faut bien en convenir : Si Dieu dans sa miséricorde avoit daigné suspendre les coups de sa justice, nous étions bien loin d'avoir employé par la pénitence le véritable moyen de fléchir sa colère. Nous devions donc nous attendre à voir de nouveau sa verge redoutable s'appesantir sur nous. Les vieux révolutionnaires s'aperçoivent en frémissant que le pouvoir va leur échapper. Ni les coups d'état ni les proscriptions n'étoient épargnés par eux pour conserver leur puissance. Les directeurs,
% Source : PDF original, vue 014, côté droit, page manuscrite 27.
des députés, sont bannis, d'autres sont expulsés des conseils ; toutes les élections qui déplaisoient au parti dominant sont cassées, et la pauvre France se voit plus que jamais enfoncée dans le gouvernement de boue du Directoire ; c'est cet événement arrivé en 7\up{bre} 1797 qu'on a appelé le 18 fructidor.

La réputation justement acquise au gouvernement de cette époque est d'avoir été souverainement méprisable ; mais il ne doit pas moins être considéré comme un gouvernement persécuteur et même sanguinaire. Les émigrés rentrés furent obligés de nouveau de quitter la France. Ceux qui avoient le malheur d'être arrêtés étoient fusillés. Les églises dont les catholiques s'étoient remis en possession furent immédiatement fermées. Il fallut apporter de nouveau la plus stricte prudence dans l'exercice de la religion, car si en général on ne guillotinoit plus les prêtres à cette époque, ce qui pourtant est encore arrivé, on les déportoit à Sinnamari, ou bien on les entassoit à Rochefort à bord de bâtiments dans lesquels ils étoient condamnés à supporter pour la foi des souffrances pires que la mort. Ce malheureux temps a été qualifié du nom assez expressif de Demi-Terreur.

On fut obligé de fermer bien vite le nouveau séminaire ; mais le zèle des catholiques dirigé par les vicaires généraux sut y suppléer. Trente maisons se chargèrent de recevoir tour à tour les étudiants pendant une journée. Ceux-ci, logés en ville par les soins de l'administration du diocèse, se rendoient chaque matin dans le lieu où se suivoit la

% Lacune constatée entre deux vues consécutives du PDF fourni.
\begin{quote}
\textit{[Lacune dans le document numérisé : la page manuscrite 27, vue PDF 014, est suivie directement de la page manuscrite 30, vue PDF 015. Les pages 28 et 29 ne figurent pas à cet emplacement. Le passage manquant n'est pas reconstitué.]}
\end{quote}

% Source : PDF original, vue 015, côté gauche, page manuscrite 30.
Rhétorique, en l'honneur de la Ste République. Vous pensez bien que les catholiques ne prenoient aucune part à toutes ces fêtes, absurdes ou ridicules, mais toujours païennes ; ils se tenoient ces jours-là enfermés dans l'intérieur de leurs maisons où ils gémissoient en secret et s'abstenoient même d'ouvrir leurs croisées pour jeter sur le cortège un coup d'œil de simple curiosité. Au surplus le bon sens de notre population apprécia bientôt à leur juste valeur ces prétendues solennités ; elle ne prenoit d'autre part que de les choisir pour but de ses sarcasmes et de ses plaisanteries.

Vous voyez, mes frères, à quel état nous étions alors réduits. Nous qui demeurions fidèles à notre religion ne pouvions que furtivement et en secret satisfaire à ses consolantes pratiques. Heureux encore nous n'étions pas entièrement privés de ses secours. Malgré toutes les recommandations envoyées de Paris par les autorités supérieures, la persécution n'étoit pas bien active dans notre ville. L'esprit de la cité étoit alors absolument opposé à toutes les vexations qu'on auroit voulu exercer sur nous. La commission de trois membres qui, sous le nom de bureau central, étoit chargée de la police de la ville n'osoit pas, à moins de dénonciations formelles, se porter à des actes de violence. Il lui étoit d'ailleurs impossible d'introduire des espions auprès des catholiques, et ceux-ci ne manquoient pas à l'occasion de s'avertir des dangers qui pouvoient les menacer. Malgré cela il étoit toujours nécessaire
% Source : PDF original, vue 015, côté droit, page manuscrite 31.
d'user de beaucoup de prudence et de se tenir constamment sur ses gardes, car les dénonciations d'un ennemi ou d'un énergumène d'impiété suffisoient pour amener une visite domiciliaire, et si elle occasionnoit l'arrestation d'un prêtre, sa déportation étoit presque inévitable. Heureusement les autorités n'encourageoient pas les dénonciateurs. Leur zèle, comme je vous le disois tout à l'heure, étoit singulièrement refroidi par l'attitude de la population qui prenoit ouvertement le parti des persécutés. Voici à ce sujet une petite scène dont j'ai été personnellement témoin. J'entrois dans la rue Buisson ; je la trouve remplie d'une foule nombreuse. Tous les yeux paroissoient se fixer avec inquiétude sur la porte d'entrée d'une maison ; je m'approche, je m'informe de ce dont il est question ; j'apprends qu'un prêtre a été dénoncé et que depuis deux heures des commissaires de police escortés de la force armée cherchent à le découvrir dans cette maison. Partageant alors l'inquiétude et la stupeur commune, je me mêle à la foule et comme les autres j'attends pour connoître le dénouement de cette triste affaire. Il s'écoule encore à peu près une demi-heure ; enfin on aperçoit les commissaires sur la porte de l'allée ; ils sortent lentement et la tête basse. À l'instant, de tous les côtés, on entend répéter dans la foule avec l'accent du bonheur : Ils ne l'ont pas trouvé ! Il n'y a personne d'arrêté ! Ils n'ont pris personne !... Ce peuple qui paroissoit il n'y a qu'un moment si consterné relève fièrement la tête et fait entendre
% Source : PDF original, vue 016, côté gauche, page manuscrite 32.
de joyeux applaudissements, et les sicaires de la persécution se hâtent de se dérober par une prompte retraite aux marques trop visibles du peu de sympathie qu'ils inspirent.

C'étoit comme vous n'en doutez pas une véritable jouissance pour nous lorsque nous pouvions être les témoins, comme dans cette circonstance, d'un petit désappointement pour nos persécuteurs ; mais cela n'empêchoit pas que pendant ces tristes années la jeunesse chrétienne ne fût bien à plaindre. Comme la prudence forçoit à se maintenir dans un état d'isolement, il devenoit presque impossible de trouver un petit cercle d'amis où on pût librement ouvrir son âme. Vous devez comprendre par ce que vous voyez encore aujourd'hui dans nos églises que les femmes, ce sexe que l'Église appelle par excellence le sexe dévot, remplissoient en grand nombre nos petits oratoires, et que les hommes et surtout les jeunes gens s'y trouvoient souvent seuls ou presque seuls. Ils étoient alors réduits à remplir au pied de la croix, au milieu des saintes femmes, le rôle de St Jean l'Évangéliste ; avec une foi plus vive nous nous fussions pénétrés des mêmes sentimens qui remplissoient sur le Calvaire le cœur de ce disciple bien-aimé. Mais la foiblesse de notre humanité avoit de la peine à se faire à ce rôle. Ajoutez encore qu'il falloit se résigner à se voir réduit pour tout exercice de religion à une messe basse, encore souvent
% Source : PDF original, vue 016, côté droit, page manuscrite 33.
célébrée par prudence au milieu de la nuit.

Aussi que de fois à cette malheureuse époque, lorsque en parcourant nos rues je considérois ces places desquelles on avoit arraché le signe du salut, ces temples tristement fermés et tombant en ruines, ces clochers réduits au silence, que de fois n'ai-je pas senti mon cœur se serrer ! Mes yeux se mouilloient de larmes et dans la douleur qui oppressoit mon âme j'élevois mes soupirs vers notre bon maître : Jusques à quand, Seigneur, lui disois-je alors, jusques à quand abandonnerez-vous votre héritage ? Jusques à quand permettrez-vous, ô mon Dieu, aux ennemis de votre saint nom de fouler sous leurs pieds notre malheureuse nation ? Jusques à quand, dans la fureur de leur impiété, feront-ils cesser les fêtes que nous célébrions en votre honneur ? Jusques à quand resterons-nous sans temples et sans autels ? Ce lien admirable de prières d'une part, de grâces et de protection de l'autre, qui unit le ciel avec la terre, ne sera-t-il donc jamais rétabli ? C'est ainsi qu'assis comme autrefois les enfants d'Israël sur les rivages des fleuves de Babylone nous exhalions nos plaintes, nous soupirions et nous gémissions avec eux en nous rappelant les beaux jours de Sion. Dieu dans sa miséricorde prêtoit une oreille attentive à notre prière. Il avoit décidé dans ses conseils éternels de venir à notre aide ; mais pour que l'action toute miséricordieuse de sa Providence fût plus sensible aux yeux de
% Source : PDF original, vue 017, côté gauche, page manuscrite 34.
tous, il falloit qu'auparavant notre foi fût soumise à de nouvelles épreuves.

Suivons donc ici la marche de cette divine Providence ; élevons nos cœurs vers Dieu, ne nous lassons pas d'admirer comment sa sagesse conduit tous les événements pour le plus grand bien des élus.

\section*{Pie VI}
% Mention marginale de la page manuscrite 34.
Nous avons déjà observé que le châtiment infligé à la France avoit été justement mérité pour l'empressement avec lequel elle avoit accueilli les doctrines mensongères d'une philosophie impie ; mais les nations voisines n'avoient que trop imité ce coupable empressement. La colère de Dieu devoit donc aussi les atteindre. Dans sa justice, le Très-Haut charge la France elle-même de châtier les peuples que les principes corrupteurs sortis de son sein avoient pervertis. Les armées de la République envahissent l'Europe. Bientôt de nouvelles républiques fabriquées à l'image de la République française s'élèvent de toutes parts ; ce sont autant de satellites qui se meuvent dans la même sphère. Ainsi nous avons fondé la République batave, la République cisalpine ; la Suisse eut beau nous répéter qu'elle étoit libre avant nous et plus libre que nous, il fallut qu'elle devînt libre à notre manière et qu'elle acceptât la République helvétique. Mais ce qui importoit par-dessus tout aux ennemis de la religion, alors maîtres du pouvoir, c'étoit de ressusciter l'antique République romaine ; car ils espéroient anéantir, écraser enfin, comme parloit Voltaire, le règne de J.-C. en détrônant celui qui le représente
% Source : PDF original, vue 017, côté droit, page manuscrite 35.
sur la terre. Les généraux de la République reçoivent donc l'ordre de marcher sur Rome. Hors d'état de se défendre, cette capitale du monde chrétien tombe aussitôt au pouvoir des révolutionnaires. Le pape Pie VI, malgré son grand âge et ses vertus, en est arraché et conduit captif jusqu'à Grenoble ; le Directoire s'étoit proposé de le faire transférer dans le centre de la France. Nous espérâmes un instant de le voir passer à Lyon ; nous eussions été si heureux de lui prouver que parmi les Lyonnais il pouvoit encore rencontrer des cœurs sincèrement attachés à la foi de leurs pères. Mais le gouvernement vit avec inquiétude l'effet prodigieux que produisoient la sainteté et la fermeté courageuse de sa victime ; sa présence suffisoit pour éveiller les sentimens de religion chez tous ceux qui avoient le bonheur de l'approcher. Pour mettre fin à ce résultat de la persécution sur lequel le Directoire n'avoit pas compté, il se hâta de faire transporter le St Pontife à Valence. À son arrivée, le peuple assemblé en foule autour de la maison où il est déposé demande à grands cris à voir le St Père et à recevoir sa bénédiction. Le St vieillard, courbé sous le poids des années, pâle de faim, accablé de fatigue, se présente à une fenêtre. Il s'identifie par sa pensée avec celui dont il est le vicaire ; il unit les outrages dont il est abreuvé et les souffrances qu'il endure avec les outrages et les souffrances auxquelles N. S. a daigné s'assujettir pour nous dans le cours de sa
% Source : PDF original, vue 018, côté gauche, page manuscrite 36.
passion, et pour faire comprendre à ce peuple fidèle le sentiment qui remplit son âme, il lui adresse cette parole sublime : \textit{Ecce homo !}\footnote{Évangile selon saint Jean 19, 5 : « Voici l'homme ! » (Note du transcripteur).} Et aussitôt, au nom de celui qu'un pareil langage désignoit ainsi, il répand la plus abondante bénédiction sur tous ceux qui l'entourent.

Peu de temps après, ce grand et saint pontife qui avoit tenu d'une main ferme au milieu d'une mer agitée le gouvernail de la barque de Pierre pendant près de 25 ans remit sa belle âme entre les mains de celui dont il avoit si fidèlement servi la cause, et fut recevoir la récompense réservée à son courage et à toutes ses vertus.

Quel moment critique pour l'Église ! Quel événement terrible que la mort de son chef dans de semblables circonstances ! Depuis longtemps l'impiété ne prenoit plus la peine de voiler ses odieux projets ; gardez bien votre pape, nous disoit-elle dans son infernale joie avec un sourire moqueur, gardez bien votre pape, car ce sera le dernier. Tout sembloit, hélas ! concourir aujourd'hui pour faire prévoir un accomplissement trop certain à cette douloureuse prophétie. L'Église, humainement parlant toute perdue, nos cœurs sont en proie à l'anxiété la plus vive. Cependant l'Évangile a dit : les portes de l'enfer ne prévaudront pas contre l'Église... L'instant décisif est donc arrivé. Qui sera vainqueur ? L'Évangile ou la philosophie ? Où est notre foi ? J.-C. ou Bélial ?... Chrétiens, que craignez-vous, hommes timides ?

% Source : PDF original, vue 018, côté droit, page manuscrite 37.
Le bras du Seigneur est-il raccourci ? N'est-ce pas lui qui a dit à la mer : Tu n'iras pas plus loin ? N'est-ce pas lui qui a cimenté sa religion par le sang des martyrs ? N'est-ce pas lui qui a fait briller aux yeux de Constantin cet étendard de la croix qui lui a donné la victoire ? Eh bien ! le même Dieu a dit aussi à son Église : Voilà que je suis avec vous jusqu'à la consommation des siècles ; fiez-vous à sa parole, il saura bien triompher de toute l'astuce des ennemis de votre foi. Vous entendez les trompettes de l'école voltairienne qui d'avance célèbrent son triomphe. Mais n'apercevez-vous pas l'Église en deuil qui vient en gémissant implorer les secours de son époux ? Seigneur, Dieu des armées, déjà votre voix se fait entendre ! Cette voix puissante qui gouverne le monde, et les légions du nord sont ébranlées ; les schismatiques ont eu la mission de sauver votre Église. Les armées révolutionnaires fuient épouvantées ! L'Italie est délivrée. Le conclave s'assemble. Le saint évêque d'Imola est placé sous le nom de Pie VII sur la chaire de St Pierre. L'impiété vaincue va cacher sa honte en frémissant de rage, et l'Église consolée s'appuyant de nouveau sur le bras protecteur qui ne lui manque jamais célèbre avec transport cette nouvelle victoire en s'écriant : Le Christ commande, le Christ règne, le Christ est vainqueur !
